% GENERATED FILE. Edit canonical lessons, then run npm run generate:books.
% canonical-source-hash: fnv1a64:b63e7ba61220920f
% canonical-lessons: ES-C564-personaje, ES-C564-protagonista, ES-C564-heroe, ES-C564-leyenda, ES-C564-mito

\chapter{El héroe de la historia}
\label{ch:es-el-heroe-de-la-historia}
\hlchaptermodality{564}

\begin{chapteropening}
\textbf{By the end of this chapter:} \emph{I can talk about the people in a story: the character, the main character, the hero, and the legends and myths they come from.}

Complete the last lesson of chapter 564: I can talk about the people in a story: the character, the main character, the hero, and the legends and myths they come from.
\end{chapteropening}

\section[el personaje]{el personaje — the character (in a story)}

\label{lesson:ES-C564-personaje}



\begin{quote}
Before the new one: say the Spanish for to relate, then the Spanish for the narrator.
\end{quote}

\subsection*{You'll want to know: el personaje}
\textbf{el personaje} — \textquotedblleft{}the character (in a story)\textquotedblright{}.

\begin{quote}
Mi personaje favorito es el abuelo.
\end{quote}

\begin{quote}
\emph{My favourite character is the grandfather.}
\end{quote}

Always masculine, even for a woman: \emph{el personaje principal}.

\begin{cousinweb}
\emph{personaje} was borrowed from French \emph{personnage} (EtymDB), built on Latin \emph{persōna} 'an actor's mask, a role'.
\end{cousinweb}

\begin{grammarlens}[title={What you've built}]
One more word for telling, explaining and describing at B1.
\end{grammarlens}

\subsection*{Your turn}
\begin{itemize}
\raggedright
  \item \emph{Say it:} \emph{el personaje}
  \item \emph{Say it:} \emph{el personaje}, once more
  \item \emph{Say it:} a sentence of your own with \emph{el personaje}
  \item \emph{Recall:} say the Spanish for to relate, then the Spanish for the narrator, then say \emph{el personaje} again
\end{itemize}

\begin{tcolorbox}[breakable,colback=teal!4,colframe=teal!35!black,title={Before you move on}]
What does el personaje mean? (\textquotedblleft{}The character (in a story)\textquotedblright{}.) Say it once more.
\end{tcolorbox}
\section[el protagonista]{el protagonista — the main character, the lead}

\label{lesson:ES-C564-protagonista}



\begin{quote}
Before the new one: say the Spanish for the narrator, then the Spanish for the character.
\end{quote}

\subsection*{You'll want to know: el protagonista}
\textbf{el protagonista} — \textquotedblleft{}the main character, the lead\textquotedblright{}.

\begin{quote}
La protagonista de la película es una médica.
\end{quote}

\begin{quote}
\emph{The main character in the film is a doctor.}
\end{quote}

The same form for a woman: \emph{la protagonista}.

\begin{cousinweb}
\emph{protagonista} is usually traced to Greek \emph{prōtagōnist\'{\={e}}s} 'the first actor', built on \emph{prôtos} 'first' and \emph{agōnist\'{\={e}}s} 'a competitor', from \emph{ag\'{\={o}}n} 'a contest'.
\end{cousinweb}

\begin{grammarlens}[title={What you've built}]
One more word for telling, explaining and describing at B1.
\end{grammarlens}

\subsection*{Your turn}
\begin{itemize}
\raggedright
  \item \emph{Say it:} \emph{el protagonista}
  \item \emph{Say it:} \emph{el protagonista}, once more
  \item \emph{Say it:} a sentence of your own with \emph{el protagonista}
  \item \emph{Recall:} say the Spanish for the narrator, then the Spanish for the character, then say \emph{el protagonista} again
\end{itemize}

\begin{tcolorbox}[breakable,colback=teal!4,colframe=teal!35!black,title={Before you move on}]
What does el protagonista mean? (\textquotedblleft{}The main character, the lead\textquotedblright{}.) Say it once more.
\end{tcolorbox}
\section[el héroe]{el héroe — the hero}

\label{lesson:ES-C564-heroe}



\begin{quote}
Before the new one: say the Spanish for the character, then the Spanish for the main character.
\end{quote}

\subsection*{You'll want to know: el héroe}
\textbf{el héroe} — \textquotedblleft{}the hero\textquotedblright{}.

\begin{quote}
El bombero fue el héroe del día.
\end{quote}

\begin{quote}
\emph{The firefighter was the hero of the day.}
\end{quote}

A woman is \emph{la heroína}.

\begin{cousinweb}
\emph{héroe} is a learned borrowing from Latin \emph{hērōs}, which comes from Greek \emph{h\'{\={e}}rōs} 'a demigod, a hero' (EtymDB).
\end{cousinweb}

\begin{grammarlens}[title={What you've built}]
One more word for telling, explaining and describing at B1.
\end{grammarlens}

\subsection*{Your turn}
\begin{itemize}
\raggedright
  \item \emph{Say it:} \emph{el héroe}
  \item \emph{Say it:} \emph{el héroe}, once more
  \item \emph{Say it:} a sentence of your own with \emph{el héroe}
  \item \emph{Recall:} say the Spanish for the character, then the Spanish for the main character, then say \emph{el héroe} again
\end{itemize}

\begin{tcolorbox}[breakable,colback=teal!4,colframe=teal!35!black,title={Before you move on}]
What does el héroe mean? (\textquotedblleft{}The hero\textquotedblright{}.) Say it once more.
\end{tcolorbox}
\section[la leyenda]{la leyenda — the legend; the caption}

\label{lesson:ES-C564-leyenda}



\begin{quote}
Before the new one: say the Spanish for the main character, then the Spanish for the hero.
\end{quote}

\subsection*{You'll want to know: la leyenda}
\textbf{la leyenda} — \textquotedblleft{}the legend; the caption\textquotedblright{}.

\begin{quote}
Según la leyenda, aquí vivía un dragón.
\end{quote}

\begin{quote}
\emph{According to the legend, a dragon lived here.}
\end{quote}

On a map or a photo it is the key or caption.

\begin{cousinweb}
\emph{leyenda} comes from Medieval Latin \emph{legenda} 'things to be read' (EtymDB), from \emph{legere} 'to read'. The first legends were the lives of saints read aloud.
\end{cousinweb}

\begin{grammarlens}[title={What you've built}]
One more word for telling, explaining and describing at B1.
\end{grammarlens}

\subsection*{Your turn}
\begin{itemize}
\raggedright
  \item \emph{Say it:} \emph{la leyenda}
  \item \emph{Say it:} \emph{la leyenda}, once more
  \item \emph{Say it:} a sentence of your own with \emph{la leyenda}
  \item \emph{Recall:} say the Spanish for the main character, then the Spanish for the hero, then say \emph{la leyenda} again
\end{itemize}

\begin{tcolorbox}[breakable,colback=teal!4,colframe=teal!35!black,title={Before you move on}]
What does la leyenda mean? (\textquotedblleft{}The legend; the caption\textquotedblright{}.) Say it once more.
\end{tcolorbox}
\section[el mito]{el mito — the myth}

\label{lesson:ES-C564-mito}



\begin{quote}
Before the new one: say the Spanish for the hero, then the Spanish for the legend.
\end{quote}

\subsection*{You'll want to know: el mito}
\textbf{el mito} — \textquotedblleft{}the myth\textquotedblright{}.

\begin{quote}
Es un mito que los toros odian el rojo.
\end{quote}

\begin{quote}
\emph{It is a myth that bulls hate red.}
\end{quote}

\begin{cousinweb}
\emph{mito} is a learned borrowing from Late Latin \emph{mythos}, taken from Greek \emph{mŷthos} 'a word, a story' (EtymDB).
\end{cousinweb}

\begin{grammarlens}[title={What you've built}]
One more word for telling, explaining and describing at B1.
\end{grammarlens}

\subsection*{Your turn}
\begin{itemize}
\raggedright
  \item \emph{Say it:} \emph{el mito}
  \item \emph{Say it:} \emph{el mito}, once more
  \item \emph{Say it:} a sentence of your own with \emph{el mito}
  \item \emph{Recall:} say the Spanish for the hero, then the Spanish for the legend, then say \emph{el mito} again
\end{itemize}

\begin{tcolorbox}[breakable,colback=teal!4,colframe=teal!35!black,title={Before you move on}]
What does el mito mean? (\textquotedblleft{}The myth\textquotedblright{}.) Say it once more.
\end{tcolorbox}
